\section{Results}
\label{sec:mars_molli_results}

This section reports the navigation performance obtained over the first
\(180~\mathrm{s}\) of the Mars Molli 32-SV simulation using GPS L1 C/A-like signals.
The receiver was operated in single-frequency code-based 3D PVT mode, with RAIM/FDE
enabled and without the fixed-height 2D surface constraint. All receiver timing and
time-to-first-fix calculations use GPS time. The truth position used for the error
calculation was the Mars surface point \(\phi = 0^\circ\), \(\lambda = 0^\circ\),
\(h = 0~\mathrm{m}\), evaluated on the Mars ellipsoid used by the receiver
error-analysis pipeline \((a = 3396190~\mathrm{m}, b = 3376200~\mathrm{m})\).

For this run, the satellite-position calculation used by RTKLIB was modified to
include a Mars orbit-matching correction inside the broadcast-orbit extrapolation
path. The correction is applied during satellite state propagation, before the
least-squares PVT solution is formed, so the receiver still uses the normal decoded
telemetry, observables, pseudorange residual computation, RAIM/FDE screening, and
RTKLIB positioning chain. The purpose of the correction is not to prescribe the user
position, but to make the satellite positions computed from the Mars broadcast-like
ephemerides consistent with the simulator/SPICE orbit frame and time-tag convention.
This removes the dominant broadcast-orbit frame-fit mismatch while preserving the
receiver-side measurement and PVT estimation problem.

Within the reported \(0\)--\(180~\mathrm{s}\) interval, the receiver produced
\(326\) valid 3D PVT epochs. Valid PVT output begins at \(17.5~\mathrm{s}\), which is
the measured TTFF, and continues through the end of the reported interval at
\(180.0~\mathrm{s}\). The number of satellites used per solution varied from 6 to 9,
and the post-fit residual log contained \(2520\) used-satellite residual samples.
The first valid 3D solution was obtained at GPS TOW \(88458.5~\mathrm{s}\). Since the
simulator satellite-data record begins at GPS TOW \(88441.0~\mathrm{s}\), the
measured time-to-first-fix was \(17.5~\mathrm{s}\).
Table~\ref{tab:mars_molli_navigation_performance} summarizes the navigation
performance metrics for this first-\(180~\mathrm{s}\) window.

\begin{table}[t]
    \centering
    \caption{Navigation performance for the first \(180~\mathrm{s}\) of the Mars Molli 32-SV 3D PVT validation run.}
    \label{tab:mars_molli_navigation_performance}
    \begin{tabular}{lll}
        \toprule
        Metric & Result & Definition \\
        \midrule
        Pseudorange residual RMS
            & \(0.402~\mathrm{m}\) mean; \(0.750~\mathrm{m}\) 95\%
            & Per-epoch RMS of RTKLIB post-fit code residuals over used SVs. \\
        3D position error
            & \(7.178~\mathrm{m}\) RMS; \(8.326~\mathrm{m}\) 95\%
            & Euclidean 3D error with respect to the Mars truth position. \\
        Vertical error
            & \(1.257~\mathrm{m}\) RMS; \(2.409~\mathrm{m}\) 95\%
            & Absolute Up/radial error with respect to the Mars truth height. \\
        Receiver clock bias stability
            & \(12.87~\mathrm{ns}\) std; \(49.46~\mathrm{ns}\) peak-to-peak
            & Variation of the solved receiver clock bias over valid PVT epochs. \\
        Time-to-first-fix
            & \(17.5~\mathrm{s}\)
            & First valid 3D PVT epoch relative to the first simulator GPS TOW. \\
        \bottomrule
    \end{tabular}
\end{table}

The post-fit pseudorange residuals remained tightly bounded during the reported
interval. The mean per-epoch residual RMS was \(0.402~\mathrm{m}\), the 95th
percentile was \(0.750~\mathrm{m}\), and the maximum observed per-epoch RMS was
\(1.108~\mathrm{m}\). This indicates that, once the satellite state representation is
made consistent with the simulator trajectory, the receiver observables are mutually
consistent at the sub-meter residual level. The residual mean over used satellites was
close to zero at each epoch, which is expected because the receiver clock bias is
estimated jointly with position and absorbs the common-mode range component.

The resulting 3D position error was \(7.178~\mathrm{m}\) RMS with a 95th-percentile
value of \(8.326~\mathrm{m}\) and a maximum of \(10.111~\mathrm{m}\). The vertical
component was smaller than the full 3D error, with a 95th-percentile absolute error of
\(2.409~\mathrm{m}\). The remaining error is therefore dominated by the horizontal
components rather than by a height bias. This distinction is important for the Mars
surface case because an apparently acceptable height solution can still mask a
longitude/east component caused by satellite geometry, broadcast-frame alignment, or
residual orbit-model mismatch.

The receiver clock estimate was stable over the valid solution interval. The standard
deviation of the solved clock bias was \(12.87~\mathrm{ns}\), equivalent to
\(3.858~\mathrm{m}\) in range, with a peak-to-peak variation of \(49.46~\mathrm{ns}\)
\((14.828~\mathrm{m})\). A linear fit to the clock-bias time history gives a drift of
\(-0.075~\mathrm{m/s}\). The absolute clock bias itself is not expected to be zero,
because the receiver clock state includes the arbitrary receiver time offset; the
reported metric is therefore the stability of the estimated clock bias during the
validated PVT interval.

Inspection of the per-SV post-fit residuals shows that the clock-bias time history is
mainly driven by small geometry-dependent redistributions of common-mode range error,
rather than by a loss of receiver timing. The largest discontinuity in the clock-bias
trace occurs at approximately \(71~\mathrm{s}\), where the solved clock changes by
\(-4.58~\mathrm{m}\) (\(-15.27~\mathrm{ns}\)) between consecutive PVT epochs. This
event coincides with a change in the used satellite set in which PRN16 is removed and
PRN12 is reintroduced. PRN16 is also the most biased low-elevation satellite in the
solution set over the reported interval, with a mean post-fit residual of
\(+0.268~\mathrm{m}\), an RMS residual of \(0.532~\mathrm{m}\), and an average
elevation of only \(3.4^\circ\). Immediately before the \(71~\mathrm{s}\) transition,
PRN16 has residuals of \(+1.6\)--\(+1.9~\mathrm{m}\), while PRN6 has simultaneous
negative residuals near \(-2.1~\mathrm{m}\). When PRN16 is removed, the least-squares
solution absorbs the changed residual geometry primarily through the receiver clock
state and the post-fit residual RMS falls from \(1.108~\mathrm{m}\) to
\(0.599~\mathrm{m}\).

Other satellites contribute smaller residual structure. PRN6 is present for the
entire interval and has one of the largest residual RMS values
\((0.529~\mathrm{m})\), making it a persistent contributor to the common-mode fit.
PRN31 is the satellite most frequently associated with the largest absolute post-fit
residual in an epoch, although its mean residual remains modest
\((-0.119~\mathrm{m})\). PRN26 and PRN31 also have low apparent elevations in the
receiver geometry, so their residuals are expected to couple more strongly into the
clock and horizontal components of the solution. These SV-specific diagnostics
indicate that the remaining clock-bias pattern is dominated by low-elevation and
changing-geometry satellites---especially PRN16 and PRN6 around the \(71~\mathrm{s}\)
transition---rather than by a systematic GPS-time error.

Overall, the orbit-matched RTKLIB propagation run demonstrates that the Mars-adapted
receiver can produce stable meter-level 3D code-based PVT once the satellite
positions extrapolated from the broadcast-like ephemerides are made consistent with
the simulator reference geometry. The result also separates two effects that were
coupled in earlier broadcast-only runs: the receiver tracking/telemetry/PVT chain is
capable of producing a well-conditioned solution, but the broadcast-orbit
interpretation must match the Mars simulator frame and time-tag conventions before
that performance can be obtained without a large orbit-model-induced positioning
bias.
